\section{Other Lean Features}\label{sec:features}

\subsection{The macro system}
A macro translates one piece of syntax into another at expansion time; it
never touches the proof kernel, so it carries zero trusted cost and zero
overhead:

\begin{lstlisting}
macro "intro_true_macro" : tactic => `(tactic| exact .intro)
example : True := by intro_true_macro

macro "print_nat_zero" : command => `(command| #check (0 : Nat))
print_nat_zero
\end{lstlisting}

\subsection{Writing custom tactics}
\texttt{elab ... : tactic} gives access to the \texttt{TacticM} monad, where
one can log messages and sequence arbitrary tactics. The example below logs
an info message and then calls \texttt{simp}:

\begin{lstlisting}
open Lean Elab Tactic

elab "report_and_simp" : tactic => do
  Lean.logInfo "report_and_simp: running simp"
  evalTactic (<- `(tactic| simp))

example (n : Nat) : n + 0 = n := by report_and_simp
\end{lstlisting}

Tactics from the shared layer
\texttt{LeanPlayground.Common.Tactics} can be \texttt{import}ed by any
subproject (this very file is a use case).

\subsection{IO and functional programming}
The pure fragment (\texttt{foldl}, list combinators) can be exercised
directly with \texttt{\#eval}; side effects are confined to the \texttt{IO}
monad:

\begin{lstlisting}
def greet (name : String) : String :=
  s!"Hello, {name}, from Lean Playground!"

def squareSum (xs : List Nat) : Nat :=
  xs.foldl (fun acc x => acc + x * x) 0

#eval squareSum [1, 2, 3, 4]   -- 30

def main : IO Unit := do
  IO.println <| greet "world"
\end{lstlisting}
